\lysh{14758}{4}

\begin{alist}
\item Στο τέλος του $1$ου μήνα θα είναι κατασκευασμένα $5$ αυτοκίνητα.
Στο τέλος του $2$ου μήνα θα είναι κατασκευασμένα $18$ αυτοκίνητα.
Στο τέλος του $3$ου μήνα θα είναι κατασκευασμένα $31$ αυτοκίνητα.
Στο τέλος του $4$ου μήνα θα είναι κατασκευασμένα $44$ αυτοκίνητα.
Στο τέλος του $5$ου μήνα θα είναι κατασκευασμένα $57$ αυτοκίνητα και στο τέλος του $6$ου μήνα
θα είναι κατασκευασμένα $70$ αυτοκίνητα.
\item Τα αυτοκίνητα που είναι κατασκευασμένα στο τέλος κάθε μήνα είναι διαδοχικοί όροι
αριθμητικής προόδου με πρώτο όρο $\alpha_1 = 5$ και διαφορά $\omega = 13$ (τα αυτοκίνητα που
κατασκευάζονται κάθε μήνα αυξάνονται σταθερά κατά $13$).
\item Τα τέσσερα πρώτα χρόνια (στο τέλος του $48$ου μήνα δηλαδή) θα έχουν κατασκευαστεί:
$\alpha_{48} = \alpha_1 + 47\omega = 5 + 47 \cdot 13 = 616$ αυτοκίνητα.
\item Ζητάμε τον φυσικό αριθμό $\nu$ για τον οποίο ισχύει:
\begin{gather*}
\alpha_{\nu} \ge 250\text{, δηλαδή} \\
5 + (\nu - 1)13 \ge 250\text{, οπότε} \\
(\nu - 1) \ge \dfrac{245}{13}\text{, δηλαδή} \\
\nu \ge \dfrac{245}{13} + 1\text{και τελικά} \\
\nu \ge 19 \dfrac{11}{13}.
\end{gather*}
Οπότε μετά από $20$ μήνες θα έχει κατασκευαστεί το $250$ο αυτοκίνητο.
\end{alist}

